% !TEX root = ../main.tex
%% Partie « Évaluations » : énoncés (les corrigés sont dans corriges-evaluations.tex)

\partieEvaluations

\planEvaluations{%
  1/DS1/otNavy/Devoir surveillé \numero1/{1,2,3},
  1/DS2/otNavy/Devoir surveillé \numero2/{4,5,6},
  1/DS3/otNavy/Devoir surveillé \numero3/{7,8,9,10},
  1/DS4/otNavy/Devoir surveillé \numero4/{11,12,13},
  1/DS5/otNavy/Devoir surveillé \numero5/{14,15,16},
  1/DS6/otNavy/Devoir surveillé \numero6/{17,18,19,20},
  2/DM1/otGreenDark/Devoir maison \numero1/{1,2,6},
  2/DM2/otGreenDark/Devoir maison \numero2/{1,7,16},
  2/DM3/otGreenDark/Devoir maison \numero3/{11,15,19},
  2/DM4/otGreenDark/Devoir maison \numero4/{6,11,17},
  3/EX1/otOrange/Examen du premier semestre/{3,4,5,7,9,10},
  3/EX2/otOrange/Examen de fin d'année/{5,11,12,13,14,16,17,19,20},
  4/C1/otRed/Type CCINP/{1,2,5,6,9,10,12,13,14,15,16,17},
  4/C2/otRed/Type Mines-Ponts/{10,12,15,18},
  4/C3/otRed/Type Centrale-Supélec/{9,13,17,19},
  4/C4/otRed/Type X-ENS/{5,11,15,18,19}}

\sectionEval{DS}{Devoirs surveillés}{Six devoirs de quatre heures, calibrés
sur la progression de l'année : chacun porte sur les trois ou quatre
chapitres qui le précèdent et comporte des exercices indépendants, notés sur
$20$ points. La calculatrice est interdite, comme dans la plupart des
classes préparatoires. Le barème figure en tête de chaque exercice et,
question par question, dans le corrigé.}
% !TEX root = ../main.tex
%% Devoir surveillé n° 1 — chapitres 1 à 3
%% Résultats vérifiés par verifications/eval_ds.py.

\sujetEval{DS1}{MPSI · Devoir surveillé}{Devoir surveillé \numero1}
  {Durée : \textbf{4 heures}\\Calculatrice interdite\\Chapitres 1 à 3}{4}
  {Exercice 1/Ensembles et applications/5 pts,Exercice 2/Sommes et récurrences/5 pts,%
   Exercice 3/Trigonométrie/4 pts,Exercice 4/Nombres complexes/6 pts}
\consigne{Les quatre exercices sont indépendants. La qualité de la rédaction
et la précision des raisonnements entrent pour une part importante dans
l'évaluation ; les résultats seront encadrés.}

\begin{exoEval}{5 points}{Ensembles et applications}
On considère $f:\R\to\R$, $x\mapsto x^{2}-2x$.
\begin{questions}
  \item Déterminer $f\bigl(\intFF{-1}{2}\bigr)$ et
    $f^{-1}\bigl(\intFF{0}{3}\bigr)$.
  \item L'application $f$ est-elle injective ? surjective ? Montrer que $f$
    induit une bijection $g$ de $\intFO{1}{+\infty}$ sur
    $\intFO{-1}{+\infty}$ et expliciter $g^{-1}$.
  \item Soient $E,F$ deux ensembles, $h:E\to F$ et $A,B\subseteq E$. Montrer
    que $h(A\cap B)\subseteq h(A)\cap h(B)$, et que l'égalité a lieu pour
    toutes parties $A,B$ lorsque $h$ est injective. Donner, avec la fonction
    $f$, un exemple où l'inclusion est stricte.
\end{questions}
\end{exoEval}

\begin{exoEval}{5 points}{Sommes et récurrences}
Soit $n\in\N^{*}$.
\begin{questions}
  \item Calculer $\displaystyle\sum_{k=1}^{n}k\cdot k!$ \emph{(écrire
    $k=(k+1)-1$)}.
  \item Montrer que $\dfrac{1}{k+1}\dbinom nk=\dfrac{1}{n+1}\dbinom{n+1}{k+1}$,
    et en déduire $\displaystyle\sum_{k=0}^{n}\frac{1}{k+1}\binom nk$.
  \item Calculer la somme double $\displaystyle\sum_{1\leqslant i\leqslant j\leqslant n}i$.
  \item Montrer par récurrence que pour tout réel $x\geqslant-1$ :
    $(1+x)^{n}\geqslant1+nx$ (inégalité de Bernoulli).
\end{questions}
\end{exoEval}

\begin{exoEval}{4 points}{Trigonométrie}
\begin{questions}
  \item Résoudre dans $\intFO{0}{2\pi}$ l'équation
    $2\cos^{2}x+3\sin x=3$.
  \item Calculer la valeur exacte de $\cos\frac\pi8$.
  \item Établir $\cos3\theta=4\cos^{3}\theta-3\cos\theta$. En déduire que
    $\cos\frac\pi9$, $\cos\frac{5\pi}9$ et $\cos\frac{7\pi}9$ sont les
    trois racines de $8t^{3}-6t-1=0$, puis calculer leur somme et leur
    produit.
\end{questions}
\end{exoEval}

\begin{exoEval}{6 points}{Nombres complexes}
\begin{questions}
  \item Résoudre dans $\Cx$ l'équation $z^{2}-(3+\mathrm{i})z+4+3\mathrm{i}=0$.
  \item Pour $z\neq-1$, on pose $Z=\dfrac{z-\mathrm{i}}{z+1}$. En écrivant
    $z=x+\mathrm{i}y$, déterminer et représenter l'ensemble des points
    d'affixe $z$ tels que $Z$ soit réel, puis tels que $Z$ soit imaginaire
    pur.
  \item Montrer que pour tout $z\in\mathbb{U}\setminus\ens{1}$,
    $\dfrac{1+z}{1-z}$ est imaginaire pur, et l'exprimer à l'aide de
    l'argument de $z$.
  \item Soit $n\geqslant2$. Résoudre dans $\Cx$ l'équation
    $(z+1)^{n}=(z-1)^{n}$. Combien a-t-elle de solutions ? Où sont-elles
    situées ?
\end{questions}
\end{exoEval}

% !TEX root = ../main.tex
%% Devoir surveillé n° 2 — chapitres 4 à 6
%% Résultats vérifiés par verifications/eval_ds.py.

\sujetEval{DS2}{MPSI · Devoir surveillé}{Devoir surveillé \numero2}
  {Durée : \textbf{4 heures}\\Calculatrice interdite\\Chapitres 4 à 6}{4}
  {Exercice 1/Fonctions usuelles/6 pts,Exercice 2/Équations différentielles/5 pts,%
   Exercice 3/Suite récurrente/5 pts,Exercice 4/{Continuité, dérivabilité}/4 pts}
\consigne{Les quatre exercices sont indépendants. On veillera à justifier
soigneusement les dérivabilités et les passages à la limite.}

\begin{exoEval}{6 points}{Fonctions usuelles}
\begin{questions}
  \item Montrer que pour tout $x\geqslant0$ :
    $x-\dfrac{x^{2}}{2}\leqslant\ln(1+x)\leqslant x$.
  \item Soit $\varphi(x)=\arctan x+\arctan\dfrac1x$ sur $\R^{*}$. Calculer
    $\varphi'$ et en déduire une expression simple de $\varphi$ sur
    $\intOO{0}{+\infty}$ et sur $\intOO{-\infty}{0}$.
  \item Étudier les variations de $h:x\mapsto x^{1/x}$ sur
    $\intOO{0}{+\infty}$. En déduire, sans calculatrice, lequel des deux
    nombres $\mathrm{e}^{\pi}$ et $\pi^{\mathrm{e}}$ est le plus grand.
\end{questions}
\end{exoEval}

\begin{exoEval}{5 points}{Équations différentielles}
\begin{questions}
  \item Résoudre sur $\R$ l'équation $(1+x^{2})\,y'+2x\,y=1$, puis
    déterminer la solution telle que $y(0)=1$. \emph{On pourra
    reconnaître une dérivée.}
  \item Résoudre sur $\R$ l'équation $y''-3y'+2y=\mathrm{e}^{x}$, puis
    déterminer la solution vérifiant $y(0)=y'(0)=0$.
\end{questions}
\end{exoEval}

\begin{exoEval}{5 points}{Suite récurrente}
On définit la suite $(u_n)$ par $u_0=0$ et $u_{n+1}=\sqrt{2+u_n}$.
\begin{questions}
  \item Montrer que $0\leqslant u_n<2$ et que $(u_n)$ est croissante. En
    déduire qu'elle converge, et déterminer sa limite.
  \item Montrer que $|u_{n+1}-2|\leqslant\frac12|u_n-2|$, puis que
    $|u_n-2|\leqslant\dfrac{2}{2^{n}}$.
  \item Montrer par récurrence que $u_n=2\cos\dfrac{\pi}{2^{n+1}}$.
  \item En déduire $\displaystyle\lim_{n\to+\infty}4^{n}(2-u_n)$.
\end{questions}
\end{exoEval}

\begin{exoEval}{4 points}{Continuité, dérivabilité}
\begin{questions}
  \item Soit $f$ continue sur $\intFF{0}{1}$ avec $f(0)=f(1)$. Montrer
    qu'il existe $c\in\intFF{0}{\frac12}$ tel que
    $f\bigl(c+\frac12\bigr)=f(c)$. \emph{Introduire
    $g(x)=f\bigl(x+\frac12\bigr)-f(x)$.} Interpréter géométriquement.
  \item Soit $P\in\R[X]$ possédant $p\geqslant2$ racines réelles distinctes.
    Montrer que $P'$ possède au moins $p-1$ racines réelles distinctes. En
    déduire que, pour $n\geqslant2$ et $a,b\in\R$, le polynôme
    $X^{n}+aX+b$ a au plus trois racines réelles distinctes.
\end{questions}
\end{exoEval}

% !TEX root = ../main.tex
%% Devoir surveillé n° 3 — chapitres 7 à 10
%% Résultats vérifiés par verifications/eval_ds.py.

\sujetEval{DS3}{MPSI · Devoir surveillé}{Devoir surveillé \numero3}
  {Durée : \textbf{4 heures}\\Calculatrice interdite\\Chapitres 7 à 10}{4}
  {Exercice 1/Arithmétique/5 pts,Exercice 2/Groupes/5 pts,%
   Exercice 3/Calcul matriciel/5 pts,Exercice 4/Polynômes et fractions/5 pts}
\consigne{Les quatre exercices sont indépendants et peuvent être traités
dans l'ordre de son choix.}

\begin{exoEval}{5 points}{Arithmétique}
\begin{questions}
  \item Résoudre dans $\Z^{2}$ l'équation $21x+15y=9$.
  \item Déterminer les deux derniers chiffres de l'écriture décimale de
    $7^{2026}$.
  \item Montrer que pour tout $n\in\Z$, $n^{5}-n$ est divisible par $30$.
  \item Pour $n\in\N$, montrer que $(n^{2}+1)\wedge(n+1)\in\ens{1,2}$, et
    préciser quand il vaut $2$.
\end{questions}
\end{exoEval}

\begin{exoEval}{5 points}{Groupes}
Soit $n\in\N^{*}$ et $\omega=\mathrm{e}^{2\mathrm{i}\pi/n}$.
\begin{questions}
  \item Montrer que $\mathbb{U}_n$ est un sous-groupe de $(\Cx^{*},\times)$,
    puis que $\varphi:k\mapsto\omega^{k}$ est un morphisme surjectif de
    $(\Z,+)$ sur $\mathbb{U}_n$ dont on déterminera le noyau.
  \item Montrer que les sous-groupes de $(\Z,+)$ sont les $m\Z$,
    $m\in\N$. \emph{Considérer le plus petit élément strictement positif.}
  \item Montrer que l'ordre de $\omega^{k}$ est $\dfrac{n}{n\wedge k}$.
    Quels éléments de $\mathbb{U}_{12}$ engendrent ce groupe ?
\end{questions}
\end{exoEval}

\begin{exoEval}{5 points}{Calcul matriciel}
\begin{questions}
  \item Soit $A=\begin{pmatrix}1&1&1\\0&1&1\\0&0&1\end{pmatrix}$. En
    écrivant $A=I_3+N$, calculer $A^{n}$ pour tout $n\in\N$.
  \item Montrer que $A$ est inversible et que la formule précédente reste
    vraie pour $n=-1$.
  \item Résoudre, selon la valeur du réel $m$, le système
    \[
    \begin{cases}x+y+mz=1\\ x+my+z=1\\ mx+y+z=1\end{cases}
    \]
\end{questions}
\end{exoEval}

\begin{exoEval}{5 points}{Polynômes et fractions}
On note $P=X^{4}+2X^{3}+3X^{2}+2X+1$ et $F=\dfrac{1}{X^{3}-X}$.
\begin{questions}
  \item Développer $(X^{2}+X+1)^{2}$, puis factoriser $P$ dans $\Cx[X]$ et
    dans $\R[X]$. Retrouver la multiplicité de la racine $j$ à l'aide des
    dérivées.
  \item Pour $n\geqslant2$, déterminer le reste de la division euclidienne
    de $X^{n}$ par $(X-1)^{2}$.
  \item Décomposer $F$ en éléments simples, et en déduire une primitive de
    $x\mapsto F(x)$ sur $\intOO{1}{+\infty}$.
\end{questions}
\end{exoEval}

% !TEX root = ../main.tex
%% Devoir surveillé n° 4 — chapitres 11 à 13
%% Résultats vérifiés par verifications/eval_ds.py.

\sujetEval{DS4}{MPSI · Devoir surveillé}{Devoir surveillé \numero4}
  {Durée : \textbf{4 heures}\\Calculatrice interdite\\Chapitres 11 à 13}{4}
  {Exercice 1/Développements limités/6 pts,Exercice 2/Suite implicite/4 pts,%
   Exercice 3/Espaces vectoriels/5 pts,Exercice 4/Projecteur et matrices/5 pts}
\consigne{Les quatre exercices sont indépendants. Les développements limités
seront présentés avec leur reste.}

\begin{exoEval}{6 points}{Développements limités}
\begin{questions}
  \item Déterminer le développement limité à l'ordre $4$ en $0$ de
    $\mathrm{e}^{\sin x}$.
  \item En déduire $\displaystyle\lim_{x\to0}\frac{\mathrm{e}^{\sin x}-\mathrm{e}^{x}}{x^{3}}$.
  \item Montrer que $\sqrt{x^{2}+2x}=x+1-\dfrac{1}{2x}+o\Bigl(\dfrac1x\Bigr)$
    en $+\infty$. En déduire l'asymptote de la courbe de
    $x\mapsto\sqrt{x^{2}+2x}$ et la position relative.
  \item Calculer
    $\displaystyle\lim_{x\to0}\Bigl(\frac{1}{\sin^{2}x}-\frac{1}{x^{2}}\Bigr)$.
\end{questions}
\end{exoEval}

\begin{exoEval}{4 points}{Suite implicite}
Pour $n\geqslant2$, on pose $f_n(x)=x^{n}+x-1$.
\begin{questions}
  \item Montrer que $f_n$ admet une unique racine $u_n$ dans $\intOO01$.
    Calculer $u_2$.
  \item Calculer $f_{n+1}(u_n)$ et en déduire que $(u_n)$ est croissante.
    Montrer que $u_n\to1$.
  \item On pose $v_n=1-u_n$. Montrer que $n\ln(1-v_n)=\ln v_n$, puis que
    $n\,v_n\sim-\ln v_n$. En déduire que $v_n\sim\dfrac{\ln n}{n}$.
\end{questions}
\end{exoEval}

\begin{exoEval}{5 points}{Espaces vectoriels}
Soit $E=\R_3[X]$ et $F=\ens{P\in E\mid P(1)=P'(1)=0}$.
\begin{questions}
  \item Montrer que $\varphi:P\mapsto\bigl(P(1),P'(1)\bigr)$ est une
    application linéaire surjective de $E$ sur $\R^{2}$. En déduire que $F$
    est un sous-espace vectoriel de $E$ et donner sa dimension.
  \item Montrer que $F=(X-1)^{2}\,\R_1[X]$ et en donner une base.
  \item Montrer que $G=\R_1[X]$ est un supplémentaire de $F$ dans $E$, et
    décomposer $X^{3}$ selon $E=F\oplus G$.
\end{questions}
\end{exoEval}

\begin{exoEval}{5 points}{Projecteur et matrices}
Dans $\R^{3}$, on note $\mathcal{P}$ le plan d'équation $x+y+z=0$ et
$D=\operatorname{Vect}(e_1)$ avec $e_1=(1,0,0)$.
\begin{questions}
  \item Montrer que $\R^{3}=\mathcal{P}\oplus D$.
  \item Soit $p$ le projecteur sur $\mathcal{P}$ parallèlement à $D$.
    Déterminer $p(x,y,z)$, puis la matrice $M$ de $p$ dans la base
    canonique. Vérifier $M^{2}=M$ et comparer $\operatorname{tr}M$ et
    $\operatorname{rg}M$.
  \item Déterminer une base de $\R^{3}$ dans laquelle la matrice de $p$ est
    $\operatorname{diag}(1,1,0)$, et la matrice de passage correspondante.
  \item Donner la matrice de la symétrie $s$ par rapport à $\mathcal{P}$
    parallèlement à $D$, et vérifier que $s^{2}=\mathrm{id}$.
\end{questions}
\end{exoEval}

% !TEX root = ../main.tex
%% Devoir surveillé n° 5 — chapitres 14 à 16
%% Résultats vérifiés par verifications/eval_ds.py.

\sujetEval{DS5}{MPSI · Devoir surveillé}{Devoir surveillé \numero5}
  {Durée : \textbf{4 heures}\\Calculatrice interdite\\Chapitres 14 à 16}{3}
  {Exercice 1/Permutations et déterminants/6 pts,Exercice 2/Intégration/8 pts,%
   Exercice 3/Dénombrement/6 pts}
\consigne{Les trois exercices sont indépendants. Dans l'exercice 3, on
justifiera chaque dénombrement par une construction explicite.}

\begin{exoEval}{6 points}{Permutations et déterminants}
\begin{questions}
  \item Soit $\sigma\in S_8$ définie par
    $\sigma=\begin{pmatrix}1&2&3&4&5&6&7&8\\3&5&7&1&2&8&4&6\end{pmatrix}$.
    Décomposer $\sigma$ en cycles à supports disjoints ; en déduire sa
    signature et son ordre.
  \item Soient $a,b\in\R$ et $D_n$ le déterminant de la matrice de taille
    $n$ dont les coefficients diagonaux valent $a$ et les autres $b$.
    Montrer que $D_n=\bigl(a+(n-1)b\bigr)(a-b)^{n-1}$. \emph{Ajouter toutes
    les colonnes à la première.}
  \item En déduire que $J_n-I_n$ est inversible pour $n\geqslant2$, où
    $J_n$ est la matrice dont tous les coefficients valent $1$.
  \item Sans le développer, montrer que
    $\begin{vmatrix}1&1&1\\a&b&c\\b+c&c+a&a+b\end{vmatrix}=0$.
\end{questions}
\end{exoEval}

\begin{exoEval}{8 points}{Intégration}
Pour $n\in\N$, on pose $I_n=\displaystyle\int_0^{1}\frac{x^{n}}{1+x^{2}}\,\mathrm{d}x$.
\begin{questions}
  \item Calculer $I_0$ et $I_1$.
  \item Montrer que $(I_n)$ est décroissante, positive, et que
    $I_n\leqslant\dfrac{1}{n+1}$.
  \item Montrer que $I_n+I_{n+2}=\dfrac{1}{n+1}$, puis que
    $\dfrac{1}{2(n+1)}\leqslant I_n\leqslant\dfrac{1}{2(n-1)}$ pour
    $n\geqslant2$. En déduire un équivalent de $I_n$.
  \item Déterminer $\displaystyle\lim_{n\to+\infty}\sum_{k=1}^{n}\frac{1}{\sqrt{n^{2}+kn}}$.
  \item À l'aide de l'inégalité de Taylor-Lagrange, montrer que :
    \[\forall x\geqslant0,\qquad
      \Bigl|\ln(1+x)-x+\frac{x^{2}}{2}\Bigr|\leqslant\frac{x^{3}}{3}.\]
\end{questions}
\end{exoEval}

\begin{exoEval}{6 points}{Dénombrement}
\begin{questions}
  \item Combien d'anagrammes le mot MISSISSIPPI possède-t-il ?
  \item On tire $5$ cartes dans un jeu de $32$ ($8$ hauteurs, $4$
    couleurs). Combien de mains contiennent un \emph{full} (trois cartes
    d'une hauteur et deux d'une autre) ? une \emph{double paire} (deux
    paires de hauteurs distinctes et une cinquième carte d'une troisième
    hauteur) ?
  \item De combien de façons peut-on distribuer $10$ bonbons identiques à
    $4$ enfants, chacun en recevant au moins un ? Et sans cette
    contrainte ?
\end{questions}
\end{exoEval}

% !TEX root = ../main.tex
%% Devoir surveillé n° 6 — chapitres 17 à 20
%% Résultats vérifiés par verifications/eval_ds.py.

\sujetEval{DS6}{MPSI · Devoir surveillé}{Devoir surveillé \numero6}
  {Durée : \textbf{4 heures}\\Calculatrice interdite\\Chapitres 17 à 20}{4}
  {Exercice 1/L'urne de Pólya/6 pts,Exercice 2/Produit scalaire/5 pts,%
   Exercice 3/Séries numériques/5 pts,Exercice 4/Deux variables/4 pts}
\consigne{Les quatre exercices sont indépendants. On rappelle que
$\sum_{n\geqslant0}\frac{1}{n!}$ converge et a pour somme $\mathrm{e}$.}

\begin{exoEval}{6 points}{L'urne de Pólya}
Une urne contient initialement une boule rouge et une boule blanche. À
chaque étape, on tire une boule au hasard, on note sa couleur, puis on la
remet dans l'urne \emph{avec une boule supplémentaire de la même couleur}.
On note $X_n$ le nombre de boules rouges après $n$ tirages ($X_0=1$) et
$R_k$ l'événement « la $k$-ième boule tirée est rouge ».
\begin{questions}
  \item Déterminer les lois de $X_1$ et de $X_2$.
  \item Montrer par récurrence que $X_n$ suit la loi uniforme sur
    $\ens{1,\dots,n+1}$.
  \item En déduire $E(X_n)$, $V(X_n)$, puis $P(R_{n+1})$.
  \item Calculer $P_{R_2}(R_1)$. Commenter.
\end{questions}
\end{exoEval}

\begin{exoEval}{5 points}{Produit scalaire}
Sur $E=\R_2[X]$, on pose
$\langle P,Q\rangle=P(0)Q(0)+P(1)Q(1)+P(2)Q(2)$.
\begin{questions}
  \item Montrer que $\langle\cdot,\cdot\rangle$ est un produit scalaire
    sur $E$.
  \item Soient $L_0,L_1,L_2$ les polynômes de Lagrange associés à
    $0,1,2$. Montrer que $(L_0,L_1,L_2)$ est une base orthonormée de $E$
    et donner les coordonnées d'un polynôme $P$ dans cette base.
  \item Déterminer le projeté orthogonal de $X^{2}$ sur $F=\R_1[X]$ et la
    distance de $X^{2}$ à $F$.
  \item Interpréter le résultat en termes de droite approchant au mieux
    les points $(0,0)$, $(1,1)$ et $(2,4)$.
\end{questions}
\end{exoEval}

\begin{exoEval}{5 points}{Séries numériques}
\begin{questions}
  \item Déterminer la nature de
    $\sum_{n\geqslant1}\Bigl(\dfrac1n-\ln\Bigl(1+\dfrac1n\Bigr)\Bigr)$.
  \item Trouver $a,b,c\in\R$ tels que $u_n=\dfrac{an^{2}+bn+c}{2^{n}}$
    vérifie $u_n-u_{n+1}=\dfrac{n^{2}}{2^{n}}$ pour tout $n$. En déduire la
    somme $\sum_{n=1}^{+\infty}\frac{n^{2}}{2^{n}}$.
  \item On note $S_N=\sum_{n=2}^{N}\ln\Bigl(1+\dfrac{(-1)^{n}}{n}\Bigr)$.
    Calculer $S_{2N+1}$ et $S_{2N}$. En déduire que la série converge et
    donner sa somme.
  \item Montrer que $\sum_{n\geqslant0}\dfrac{n+1}{n!}$ converge et calculer
    sa somme.
\end{questions}
\end{exoEval}

\begin{exoEval}{4 points}{Deux variables}
Soit $f:(x,y)\mapsto x^{2}+xy+y^{2}-3x-3y$, définie sur $\R^{2}$.
\begin{questions}
  \item Justifier que $f$ est de classe $\mathcal{C}^{1}$, calculer ses
    dérivées partielles et déterminer ses points critiques.
  \item On pose $u=x-1$ et $v=y-1$. Montrer que
    \[f(x,y)=u^{2}+uv+v^{2}-3\qquad\text{et}\qquad
      u^{2}+uv+v^{2}=\Bigl(u+\frac v2\Bigr)^{2}+\frac34\,v^{2}.\]
    Conclure sur les extremums de $f$.
  \item Donner l'équation du plan tangent à la surface $z=f(x,y)$ en
    l'origine, et la direction de plus forte décroissance de $f$ en
    $(0,0)$.
\end{questions}
\end{exoEval}


\sectionEval{DM}{Devoirs maison}{Quatre problèmes longs, à chercher sur une
semaine. Ils prolongent le cours et préparent aux grands thèmes des épreuves
de concours ; on peut s'aider du cours, mais on rédigera chaque réponse
complètement.}
% !TEX root = ../main.tex
%% Devoir maison n° 1 — L'inégalité arithmético-géométrique (chapitres 1, 2, 6)
%% Résultats vérifiés par verifications/eval_dm.py.

\sujetEval{DM1}{MPSI · Devoir maison}{Devoir maison \numero1}
  {À rendre en une semaine\\L'inégalité arithmético-géométrique\\Chapitres 1, 2 et 6}{3}
  {Partie A/La preuve de Cauchy/Q1 à Q4,Partie B/La preuve par la tangente/Q5 et Q6,%
   Partie C/Applications/Q7 à Q11}
\enParties
\consigne{Les parties B et C peuvent être abordées indépendamment de la
partie A, en admettant l'inégalité.}

Pour $n\geqslant1$ et $x_1,\dots,x_n$ réels strictement positifs, on note
\[A_n=\frac{x_1+x_2+\dots+x_n}{n}\qquad\text{et}\qquad G_n=\bigl(x_1x_2\cdots x_n\bigr)^{1/n}\]
leurs moyennes \emph{arithmétique} et \emph{géométrique}. L'objet du problème
est de démontrer de deux façons l'inégalité $G_n\leqslant A_n$, avec égalité si
et seulement si les $x_i$ sont tous égaux, puis d'en tirer quelques
conséquences.

\begin{exoEval}[La preuve de Cauchy]{}{}
Pour $n\geqslant1$, on note $\mathcal{P}(n)$ l'assertion : « pour tous réels
$x_1,\dots,x_n>0$, on a $G_n\leqslant A_n$ ».
\begin{questions}
  \item Démontrer $\mathcal{P}(2)$ et préciser le cas d'égalité.
  \item Soit $n\geqslant1$ tel que $\mathcal{P}(n)$ soit vraie, et
    $x_1,\dots,x_{2n}>0$. On note $A',G'$ (resp. $A'',G''$) les moyennes
    arithmétique et géométrique de $x_1,\dots,x_n$ (resp. de
    $x_{n+1},\dots,x_{2n}$). Montrer que $G_{2n}=\sqrt{G'G''}$, puis en
    déduire $\mathcal{P}(2n)$. Conclure que $\mathcal{P}(2^{k})$ est vraie
    pour tout $k\in\N$.
  \item Soit $n\geqslant2$ tel que $\mathcal{P}(n)$ soit vraie, et
    $x_1,\dots,x_{n-1}>0$ de moyenne arithmétique $A$. En appliquant
    $\mathcal{P}(n)$ aux $n$ réels $x_1,\dots,x_{n-1},A$, montrer
    $\mathcal{P}(n-1)$.
  \item En déduire que $\mathcal{P}(n)$ est vraie pour tout $n\geqslant1$.
    \emph{On parle de récurrence « en avant, puis en arrière ».}
\end{questions}
\end{exoEval}

\begin{exoEval}[La preuve par la tangente]{}{}
\begin{questions}[start=5]
  \item Montrer que pour tout $t>0$, $\ln t\leqslant t-1$, avec égalité si
    et seulement si $t=1$. Interpréter graphiquement.
  \item En appliquant cette inégalité aux réels $t_i=\dfrac{x_i}{A_n}$,
    retrouver l'inégalité $G_n\leqslant A_n$ et montrer qu'il y a égalité si
    et seulement si $x_1=x_2=\dots=x_n$.
\end{questions}
\end{exoEval}

\begin{exoEval}[Applications]{}{}
\begin{questions}[start=7]
  \item On appelle moyenne \emph{harmonique} des $x_i$ le réel
    $H_n=\dfrac{n}{\frac{1}{x_1}+\dots+\frac{1}{x_n}}$. Montrer que
    $H_n\leqslant G_n\leqslant A_n$.
  \item Soit $n\geqslant1$. En appliquant l'inégalité aux $n+1$ réels
    $1+\frac1n,\dots,1+\frac1n$ ($n$ fois) et $1$, montrer que la suite de
    terme général $u_n=\bigl(1+\frac1n\bigr)^{n}$ est strictement
    croissante.
  \item Montrer que $n!\leqslant\Bigl(\dfrac{n+1}{2}\Bigr)^{n}$ pour tout
    $n\geqslant1$, avec égalité seulement pour $n=1$.
  \item En écrivant $x=\frac x2+\frac x2$, déterminer le minimum sur
    $\intOO{0}{+\infty}$ de la fonction $x\mapsto x+\dfrac{4}{x^{2}}$, sans
    la dériver.
  \item Parmi les pavés droits de dimensions $x,y,z>0$ telles que
    $x+y+z=3$, lesquels ont le plus grand volume ?
\end{questions}
\end{exoEval}

% !TEX root = ../main.tex
%% Devoir maison n° 2 — Résidus quadratiques (chapitres 1, 7, 16)
%% Résultats vérifiés par verifications/eval_dm.py.

\sujetEval{DM2}{MPSI · Devoir maison}{Devoir maison \numero2}
  {À rendre en une semaine\\Résidus quadratiques\\Chapitres 1, 7 et 16}{3}
  {Partie A/Compter les carrés/Q1 à Q3,Partie B/Les théorèmes de Wilson et d'Euler/Q4 à Q6,%
   Partie C/Applications/Q7 à Q10}
\enParties
\consigne{Dans tout le problème, $p$ désigne un nombre premier impair, et
les congruences sont prises modulo $p$ sauf mention contraire. Seuls le
lemme d'Euclide, le théorème de Bézout et le petit théorème de Fermat sont
supposés connus.}

Un entier $a$ non divisible par $p$ est appelé \emph{résidu quadratique}
modulo $p$ s'il existe $x\in\Z$ tel que $x^{2}\equiv a\pmod p$, et
\emph{non-résidu} sinon. Par exemple, $2$ est un résidu modulo $7$ car
$3^{2}=9\equiv2\pmod7$. On note $R_p$ l'ensemble des résidus quadratiques
appartenant à $\ens{1,\dots,p-1}$.

\begin{exoEval}[Compter les carrés]{}{}
\begin{questions}
  \item Soient $x,y\in\ens{1,\dots,p-1}$. Montrer que
    $x^{2}\equiv y^{2}\pmod p$ si et seulement si $x=y$ ou $x+y=p$.
  \item En déduire que $R_p$ possède exactement $\dfrac{p-1}{2}$ éléments :
    les restes des divisions par $p$ de $1^{2},2^{2},\dots,\bigl(\frac{p-1}{2}\bigr)^{2}$.
  \item Déterminer $R_{13}$.
\end{questions}
\end{exoEval}

\begin{exoEval}[Les théorèmes de Wilson et d'Euler]{}{}
\begin{questions}[start=4]
  \item \begin{enumerate}[label=\alph*)]
      \item Montrer que pour tout $x\in\ens{1,\dots,p-1}$, il existe un
        unique $x^{*}\in\ens{1,\dots,p-1}$ tel que $xx^{*}\equiv1$.
      \item Montrer que $x^{*}=x$ si et seulement si $x\in\ens{1,p-1}$.
      \item En regroupant chaque $x$ avec $x^{*}$, démontrer le
        \emph{théorème de Wilson} : $(p-1)!\equiv-1\pmod p$.
    \end{enumerate}
  \item Soit $a$ un non-résidu. Montrer que pour tout
    $x\in\ens{1,\dots,p-1}$, il existe un unique $x'\in\ens{1,\dots,p-1}$
    tel que $xx'\equiv a$, et que $x'\neq x$. En regroupant les éléments de
    $\ens{1,\dots,p-1}$ par paires $\ens{x,x'}$, montrer que
    $(p-1)!\equiv a^{(p-1)/2}$, puis que $a^{(p-1)/2}\equiv-1$.
  \item Soit $a$ un résidu. Montrer que $a^{(p-1)/2}\equiv1$. En déduire le
    \emph{critère d'Euler} : pour $a$ non divisible par $p$,
    \[a\ \text{est un résidu quadratique modulo }p\iff a^{(p-1)/2}\equiv1\pmod p.\]
\end{questions}
\end{exoEval}

\begin{exoEval}[Applications]{}{}
\begin{questions}[start=7]
  \item À l'aide du critère d'Euler, déterminer si $2$ et $3$ sont des
    résidus quadratiques modulo $13$.
  \item Montrer que $-1$ est un résidu quadratique modulo $p$ si et
    seulement si $p\equiv1\pmod4$.
  \item On suppose $p\equiv1\pmod4$. En écrivant
    $(p-1)!=\prod_{k=1}^{(p-1)/2}k\,(p-k)$, montrer que
    $x=\bigl(\frac{p-1}{2}\bigr)!$ vérifie $x^{2}\equiv-1\pmod p$. Vérifier
    sur $p=13$.
  \item Montrer qu'il existe une infinité de nombres premiers congrus à
    $1$ modulo $4$. \emph{Raisonner par l'absurde en considérant
    $N=(2p_1p_2\cdots p_r)^{2}+1$.}
\end{questions}
\end{exoEval}

% !TEX root = ../main.tex
%% Devoir maison n° 3 — Méthode des trapèzes et formule de Stirling (chapitres 11, 15, 19)
%% Résultats vérifiés par verifications/eval_dm.py.

\sujetEval{DM3}{MPSI · Devoir maison}{Devoir maison \numero3}
  {À rendre en une semaine\\La formule de Stirling\\Chapitres 11, 15 et 19}{3}
  {Partie A/L'erreur de la méthode des trapèzes/Q1 à Q3,Partie B/Un développement de $\ln n!$/Q4 à Q7,%
   Partie C/La constante de Stirling/Q8 et Q9}
\enParties
\consigne{La partie C utilise le résultat de la question 6, que l'on
pourra admettre.}

La \emph{méthode des trapèzes} approche l'intégrale d'une fonction sur un
segment par l'aire du trapèze construit sur ses valeurs aux extrémités. On
en étudie l'erreur, puis on l'applique au logarithme pour établir la
\emph{formule de Stirling} :
\[n!\underset{n\to+\infty}{\sim}\sqrt{2\pi n}\,\Bigl(\frac{n}{\mathrm{e}}\Bigr)^{n}.\]

\begin{exoEval}[L'erreur de la méthode des trapèzes]{}{}
Soient $a<b$ et $f$ une fonction de classe $\mathcal{C}^{2}$ sur
$\intFF ab$. On pose
\[T(f)=(b-a)\,\frac{f(a)+f(b)}{2}\qquad\text{et}\qquad\varphi(t)=(t-a)(b-t).\]
\begin{questions}
  \item \begin{enumerate}[label=\alph*)]
      \item En intégrant deux fois par parties $\int_a^{b}\varphi f''$,
        montrer que
        \[\int_a^{b}f(t)\,\mathrm{d}t-T(f)=-\frac12\int_a^{b}\varphi(t)\,f''(t)\,\mathrm{d}t.\]
      \item Vérifier cette formule pour $f:t\mapsto t^{2}$ sur $\intFF01$.
    \end{enumerate}
  \item On suppose que $|f''|\leqslant M$ sur $\intFF ab$. Calculer
    $\int_a^{b}\varphi$ et en déduire que
    \[\Bigl|\int_a^{b}f(t)\,\mathrm{d}t-T(f)\Bigr|\leqslant\frac{M(b-a)^{3}}{12}.\]
  \item On découpe $\intFF01$ en $n$ segments de même longueur et l'on
    applique la méthode sur chacun d'eux pour approcher
    \[\int_0^{1}\frac{\mathrm{d}t}{1+t^{2}}=\frac\pi4.\]
    Montrer que la dérivée seconde de $t\mapsto\frac{1}{1+t^{2}}$ est
    majorée en valeur absolue par $2$ sur $\intFF01$, puis que l'erreur
    commise est au plus $\dfrac{1}{6n^{2}}$. Combien de segments suffisent
    pour garantir une erreur inférieure à $10^{-6}$ ?
\end{questions}
\end{exoEval}

\begin{exoEval}[Un développement de $\ln n!$]{}{}
Pour $k\geqslant1$, on pose
\[e_k=\int_k^{k+1}\ln t\,\mathrm{d}t-\frac{\ln k+\ln(k+1)}{2}.\]
\begin{questions}[start=4]
  \item Montrer que $0\leqslant e_k\leqslant\dfrac{1}{12k^{2}}$, et
    interpréter $e_k$ comme une aire.
  \item Montrer que pour tout $n\geqslant2$ :
    \[\ln n!=n\ln n-n+\frac12\ln n+1-\sum_{k=1}^{n-1}e_k.\]
  \item En déduire que la série $\sum e_k$ converge, puis qu'il existe une
    constante $C$ telle que
    $\ln n!=n\ln n-n+\frac12\ln n+C+o(1)$. En déduire qu'il existe $K>0$ tel
    que $n!\sim K\sqrt n\,\bigl(\frac n{\mathrm{e}}\bigr)^{n}$.
  \item Montrer que pour $n\geqslant2$,
    $0\leqslant\sum_{k=n}^{+\infty}e_k\leqslant\dfrac{1}{12(n-1)}$, et en
    déduire un encadrement de l'erreur dans le développement de la
    question 6.
\end{questions}
\end{exoEval}

\begin{exoEval}[La constante de Stirling]{}{}
Pour $n\in\N$, on considère les \emph{intégrales de Wallis}
\[W_n=\int_0^{\pi/2}\sin^{n}t\,\mathrm{d}t.\]
\begin{questions}[start=8]
  \item \begin{enumerate}[label=\alph*)]
      \item Montrer que $W_{n+2}=\dfrac{n+1}{n+2}\,W_n$, et en déduire
        $W_{2p}=\dfrac{(2p)!}{4^{p}\,(p!)^{2}}\cdot\dfrac\pi2$.
      \item Montrer que $(W_n)$ est décroissante et strictement positive,
        que $(n+1)W_{n+1}W_n=\dfrac\pi2$, puis que
        $W_n\sim\sqrt{\dfrac{\pi}{2n}}$.
    \end{enumerate}
  \item En calculant de deux façons un équivalent de $W_{2p}$, montrer que
    $K=\sqrt{2\pi}$. En déduire la formule de Stirling, et une valeur
    approchée de $C$.
\end{questions}
\end{exoEval}

% !TEX root = ../main.tex
%% Devoir maison n° 4 — Inégalités de concentration (chapitres 6, 11, 17)
%% Résultats vérifiés par verifications/eval_dm.py.

\sujetEval{DM4}{MPSI · Devoir maison}{Devoir maison \numero4}
  {À rendre en une semaine\\Inégalités de concentration\\Chapitres 6, 11 et 17}{3}
  {Partie A/L'inégalité de Bienaymé-Tchebychev/Q1 et Q2,Partie B/L'inégalité de Hoeffding/Q3 à Q6,%
   Partie C/Application à un sondage/Q7 à Q9}
\enParties
\consigne{On rappelle que $\operatorname{ch}t=\frac{\mathrm{e}^{t}+\mathrm{e}^{-t}}{2}$
et $\operatorname{th}t=\frac{\operatorname{sh}t}{\operatorname{ch}t}$. Les
valeurs numériques demandées en partie C se calculent à la main, avec
$\ln40\approx3{,}689$.}

Soit $n\geqslant1$ et $X_1,\dots,X_n$ des variables aléatoires
indépendantes, définies sur un même espace probabilisé fini, de même loi
\[P(X_i=1)=P(X_i=-1)=\frac12.\]
On pose $S_n=X_1+\dots+X_n$ : c'est la position, après $n$ pas, d'un
marcheur qui avance ou recule d'un pas au hasard. On cherche à majorer la
probabilité que $S_n$ s'écarte beaucoup de $0$.

\begin{exoEval}[L'inégalité de Bienaymé-Tchebychev]{}{}
\begin{questions}
  \item Calculer $E(S_n)$ et $V(S_n)$.
  \item En déduire que pour tout $a>0$ :
    $P\bigl(|S_n|\geqslant a\bigr)\leqslant\dfrac{n}{a^{2}}$.
\end{questions}
\end{exoEval}

\begin{exoEval}[L'inégalité de Hoeffding]{}{}
Soit $a>0$.
\begin{questions}[start=3]
  \item Soit $t>0$. À l'aide de l'inégalité de Markov appliquée à
    $\mathrm{e}^{tS_n}$, montrer que
    $P(S_n\geqslant a)\leqslant\mathrm{e}^{-ta}\,E\bigl(\mathrm{e}^{tS_n}\bigr)$.
  \item Montrer que $E\bigl(\mathrm{e}^{tS_n}\bigr)=(\operatorname{ch}t)^{n}$.
  \item Étudier sur $\R_{+}$ la fonction
    $g:t\mapsto\frac{t^{2}}{2}-\ln(\operatorname{ch}t)$, et en déduire que
    $\operatorname{ch}t\leqslant\mathrm{e}^{t^{2}/2}$ pour tout réel $t$.
  \item En déduire que $P(S_n\geqslant a)\leqslant\exp\bigl(-ta+\frac{nt^{2}}{2}\bigr)$
    pour tout $t>0$, puis, en choisissant $t$ au mieux,
    \[P(S_n\geqslant a)\leqslant\exp\Bigl(-\frac{a^{2}}{2n}\Bigr)\qquad\text{et}\qquad
      P\bigl(|S_n|\geqslant a\bigr)\leqslant2\exp\Bigl(-\frac{a^{2}}{2n}\Bigr).\]
\end{questions}
\end{exoEval}

\begin{exoEval}[Application à un sondage]{}{}
On lance $n$ fois une pièce équilibrée, $X_i$ valant $1$ si le $i$-ième
lancer donne « pile » et $-1$ sinon. On note $F_n$ la fréquence de « pile »
et l'on fixe $\varepsilon>0$.
\begin{questions}[start=7]
  \item Montrer que $F_n=\dfrac{n+S_n}{2n}$, puis que
    \[P\Bigl(\Bigl|F_n-\frac12\Bigr|\geqslant\varepsilon\Bigr)\leqslant\frac{1}{4n\varepsilon^{2}}
      \qquad\text{et}\qquad
      P\Bigl(\Bigl|F_n-\frac12\Bigr|\geqslant\varepsilon\Bigr)\leqslant2\,\mathrm{e}^{-2n\varepsilon^{2}}.\]
  \item On prend $\varepsilon=0{,}05$. Pour chacune des deux majorations,
    déterminer le plus petit $n$ qui garantit
    $P\bigl(|F_n-\frac12|\geqslant0{,}05\bigr)\leqslant0{,}05$.
  \item Comparer les deux majorations lorsque $n\to+\infty$. Commenter.
\end{questions}
\end{exoEval}


\sectionEval{EX}{Examens}{Deux concours blancs de quatre heures : le premier
porte sur le premier semestre (chapitres 1 à 10), le second sur l'ensemble du
programme. Ils mêlent exercices courts et problème, comme les épreuves
écrites.}
% !TEX root = ../main.tex
%% Examen du premier semestre — chapitres 1 à 10
%% Résultats vérifiés par verifications/eval_ex.py.

\sujetEval{EX1}{MPSI · Concours blanc}{Examen du premier semestre}
  {Durée : \textbf{4 heures}\\Calculatrice interdite\\Chapitres 1 à 10}{4}
  {Exercice 1/Le pentagone régulier/4 pts,Exercice 2/Une équation singulière/4 pts,%
   Exercice 3/Nombres de Mersenne/4 pts,Problème/Fibonacci et matrices/8 pts}
\consigne{Les trois exercices et le problème sont indépendants. La
présentation, la lisibilité et la précision des raisonnements entrent pour
une part importante dans l'appréciation des copies.}

\begin{exoEval}{4 points}{Le pentagone régulier}
On note $\omega=\mathrm{e}^{2\mathrm{i}\pi/5}$, $\alpha=\omega+\omega^{4}$ et
$\beta=\omega^{2}+\omega^{3}$.
\begin{questions}
  \item Montrer que $1+\omega+\omega^{2}+\omega^{3}+\omega^{4}=0$.
  \item Montrer que $\alpha=2\cos\frac{2\pi}{5}$ et $\beta=2\cos\frac{4\pi}{5}$,
    puis calculer $\alpha+\beta$ et $\alpha\beta$.
  \item En déduire les valeurs exactes de $\cos\frac{2\pi}{5}$ et de
    $\cos\frac{4\pi}{5}$.
  \item Calculer la longueur $c$ du côté du pentagone régulier inscrit dans
    le cercle unité, de sommets $1,\omega,\dots,\omega^{4}$.
\end{questions}
\end{exoEval}

\begin{exoEval}{4 points}{Une équation singulière}
On considère l'équation différentielle $(E)$ : $x\,y'-2y=x^{3}$.
\begin{questions}
  \item Résoudre $(E)$ sur $\intOO{0}{+\infty}$ et sur $\intOO{-\infty}{0}$.
  \item Déterminer toutes les fonctions $y$ dérivables sur $\R$ qui
    vérifient $(E)$ en tout point de $\R$. L'ensemble obtenu est-il une
    droite affine de solutions, comme pour une équation $y'+a(x)y=b(x)$ ?
  \item Montrer que par tout point $(x_0,y_0)$ avec $x_0\neq0$ passe une
    infinité de solutions sur $\R$, et qu'aucune solution ne passe par
    $(0,y_0)$ si $y_0\neq0$.
\end{questions}
\end{exoEval}

\begin{exoEval}{4 points}{Nombres de Mersenne}
\begin{questions}
  \item Soient $a,b\in\N^{*}$. Montrer que $2^{a}-1$ divise $2^{ab}-1$.
  \item Soient $a,b\in\N^{*}$ et $a=bq+r$ la division euclidienne de $a$
    par $b$. Montrer que $2^{a}-1=2^{r}\bigl(2^{bq}-1\bigr)+2^{r}-1$, puis que
    \[(2^{a}-1)\wedge(2^{b}-1)=2^{a\wedge b}-1.\]
  \item Montrer que si $2^{n}-1$ est premier ($n\geqslant2$), alors $n$ est
    premier. La réciproque est-elle vraie ? \emph{On pourra examiner
    $n=11$.}
\end{questions}
\end{exoEval}

\begin{exoEval}[Problème]{8 points}{Fibonacci, matrices et polynômes}
La suite de Fibonacci est définie par $F_0=0$, $F_1=1$ et
$F_{n+2}=F_{n+1}+F_n$ pour $n\in\N$. On pose
\[A=\begin{pmatrix}1&1\\1&0\end{pmatrix},\qquad
  \varphi=\frac{1+\sqrt5}{2},\qquad\psi=\frac{1-\sqrt5}{2}.\]
\partieExo{A}{Puissances de $A$}
\begin{questions}
  \item Montrer que pour tout $n\geqslant1$,
    $A^{n}=\begin{pmatrix}F_{n+1}&F_n\\F_n&F_{n-1}\end{pmatrix}$.
  \item En écrivant $A^{m+n}=A^{m}A^{n}$, montrer que pour $m\geqslant0$ et
    $n\geqslant1$ : $F_{m+n}=F_{m+1}F_n+F_mF_{n-1}$. En déduire
    $F_{2n}=F_n\,(F_{n+1}+F_{n-1})$.
\end{questions}
\partieExo{B}{Un détour par les polynômes}
\begin{questions}[start=3]
  \item Montrer que pour tout $n\geqslant1$, le reste de la division
    euclidienne de $X^{n}$ par $X^{2}-X-1$ est $F_nX+F_{n-1}$.
  \item Vérifier que $A^{2}=A+I_2$, et en déduire que
    $A^{n}=F_nA+F_{n-1}I_2$ pour $n\geqslant1$.
  \item En évaluant en $\varphi$ et en $\psi$, établir la formule de Binet :
    $F_n=\dfrac{\varphi^{n}-\psi^{n}}{\sqrt5}$. En déduire un équivalent de
    $F_n$ et la limite de $\dfrac{F_{n+1}}{F_n}$.
\end{questions}
\partieExo{C}{Arithmétique}
\begin{questions}[start=6]
  \item Montrer par récurrence l'identité de Cassini :
    $F_{n+1}F_{n-1}-F_n^{2}=(-1)^{n}$ pour $n\geqslant1$.
  \item En déduire que deux termes consécutifs de la suite de Fibonacci sont
    premiers entre eux.
\end{questions}
\end{exoEval}

% !TEX root = ../main.tex
%% Examen de fin d'année — ensemble du programme
%% Résultats vérifiés par verifications/eval_ex.py.

\sujetEval{EX2}{MPSI · Concours blanc}{Examen de fin d'année}
  {Durée : \textbf{4 heures}\\Calculatrice interdite\\Tout le programme}{4}
  {Exercice 1/Permutations aléatoires/4 pts,Exercice 2/Matrices de rang 1/5 pts,%
   Exercice 3/Le folium de Descartes/4 pts,Problème/La suite $u_{n+1}=\sin u_n$/7 pts}
\consigne{Les trois exercices et le problème sont indépendants. Dans le
problème, la question 3 peut être admise pour traiter la suite.}

\begin{exoEval}{4 points}{Points fixes d'une permutation aléatoire}
Soit $n\geqslant2$. On choisit une permutation $\sigma$ de
$\ens{1,\dots,n}$ selon la loi uniforme sur $S_n$, et l'on note $X$ son
nombre de points fixes. Pour $1\leqslant i\leqslant n$, $X_i$ vaut $1$ si
$\sigma(i)=i$ et $0$ sinon.
\begin{questions}
  \item Calculer $P(X_i=1)$, puis $E(X)$.
  \item Pour $i\neq j$, calculer $P(X_i=1,\,X_j=1)$. En déduire $E(X^{2})$,
    puis $V(X)$.
  \item Calculer $P(X=n)$ et $P(X=n-1)$. À l'aide de l'inégalité de
    Bienaymé-Tchebychev, montrer que $P(X\geqslant3)\leqslant\frac14$.
  \item Déterminer la loi de $X$ pour $n=4$ et vérifier les résultats
    précédents.
\end{questions}
\end{exoEval}

\begin{exoEval}{5 points}{Matrices de rang 1}
Soient $n\geqslant2$, $U,V\in\mathcal{M}_{n,1}(\R)$ deux colonnes non nulles
et $M=U\,V^{\mathsf{T}}\in\mathcal{M}_n(\R)$.
\begin{questions}
  \item Montrer que $M$ est de rang $1$, que
    $\operatorname{tr}M=V^{\mathsf{T}}U$, et que
    $M^{2}=(\operatorname{tr}M)\,M$.
  \item En déduire que si $\operatorname{tr}M\neq0$,
    $\frac{1}{\operatorname{tr}M}M$ est la matrice d'un projecteur, et que si
    $\operatorname{tr}M=0$, $M$ est nilpotente.
  \item Montrer que $I_n+M$ est inversible si et seulement si
    $\operatorname{tr}M\neq-1$, et qu'alors
    $(I_n+M)^{-1}=I_n-\dfrac{1}{1+\operatorname{tr}M}\,M$.
  \item En utilisant une matrice semblable à $M$, montrer que
    $\det(I_n+M)=1+\operatorname{tr}M$.
  \item Application : $U=\begin{pmatrix}1\\2\\3\end{pmatrix}$ et
    $V=\begin{pmatrix}1\\1\\1\end{pmatrix}$. Calculer $(I_3+M)^{-1}$ et
    $\det(I_3+M)$.
\end{questions}
\end{exoEval}

\begin{exoEval}{4 points}{Le folium de Descartes}
Soit $f:(x,y)\mapsto x^{3}+y^{3}-3xy$, définie sur $\R^{2}$.
\begin{questions}
  \item Déterminer les points critiques de $f$.
  \item Montrer que $f$ n'admet pas d'extremum local en $(0,0)$.
  \item Montrer que
    $f(1+h,1+k)=-1+3\bigl(h^{2}-hk+k^{2}\bigr)+h^{3}+k^{3}$, puis que
    $h^{2}-hk+k^{2}\geqslant\frac12(h^{2}+k^{2})$. En déduire que $f$ admet en
    $(1,1)$ un minimum local. Est-il global ?
  \item Vérifier que la courbe $f(x,y)=0$ (le \emph{folium de Descartes})
    est paramétrée par $x=\dfrac{3t}{1+t^{3}}$, $y=\dfrac{3t^{2}}{1+t^{3}}$,
    $t\neq-1$.
\end{questions}
\end{exoEval}

\begin{exoEval}[Problème]{7 points}{La suite $u_{n+1}=\sin u_n$}
Soit $u_0\in\bigl]0,\frac\pi2\bigr]$ et $(u_n)$ définie par
$u_{n+1}=\sin u_n$.
\begin{questions}
  \item Montrer que pour tout $n$, $u_n\in\bigl]0,\frac\pi2\bigr]$, que
    $(u_n)$ est décroissante, et qu'elle converge vers $0$.
  \item Montrer que $\dfrac{1}{\sin^{2}x}-\dfrac{1}{x^{2}}\xrightarrow[x\to0]{}\dfrac13$,
    et en déduire la limite de $v_n=\dfrac{1}{u_{n+1}^{2}}-\dfrac{1}{u_n^{2}}$.
  \item (\emph{Lemme de Cesàro}) Soit $(a_n)$ une suite réelle de limite
    $\ell$. Montrer que $\dfrac{a_0+a_1+\dots+a_{n-1}}{n}\xrightarrow[n\to+\infty]{}\ell$.
  \item En déduire que $\dfrac{1}{n\,u_n^{2}}\to\dfrac13$, puis que
    $u_n\sim\sqrt{\dfrac3n}$.
  \item Déterminer la nature des séries $\sum u_n$, $\sum u_n^{2}$ et
    $\sum u_n^{3}$.
\end{questions}
\end{exoEval}


\sectionEval{CC}{Sujets type concours}{Quatre sujets inédits, conçus dans
l'esprit des épreuves écrites des grandes banques de concours, de difficulté
croissante : un sujet type CCINP, un problème type Mines-Ponts, un problème de
modélisation type Centrale-Supélec et un problème type X-ENS. Ce ne sont pas
des annales : ils ont été rédigés pour ce livre, sur le programme de
première année.}
% !TEX root = ../main.tex
%% Sujet type CCINP — quatre exercices (algèbre linéaire, probabilités,
%% polynômes, analyse). Résultats vérifiés par verifications/eval_cc.py.

\sujetEval{C1}{Sujet type CCINP · Mathématiques}{Sujet type CCINP}
  {Durée : \textbf{4 heures}\\Calculatrice interdite\\Programme de première année}{4}
  {Exercice 1/La matrice $J$/5 pts,Exercice 2/Records/5 pts,%
   Exercice 3/Polynômes de Bernoulli/5 pts,Exercice 4/L'équation de Cauchy/5 pts}
\consigne{Sujet inédit, conçu dans l'esprit de l'épreuve de mathématiques
de la banque CCINP : quatre exercices indépendants, de longueur et de
difficulté comparables. Les candidats sont invités à encadrer les résultats
de leurs calculs.}

\begin{exoEval}{5 points}{La matrice $J$}
Soit $n\geqslant2$ et $J\in\mathcal{M}_n(\R)$ la matrice dont tous les
coefficients valent $1$. Pour $a,b\in\R$, on pose $M(a,b)=aI_n+bJ$.
\begin{questions}
  \item Calculer $J^{2}$. Déterminer le rang de $J$, puis une base de son
    image et la dimension de son noyau.
  \item Montrer que $\R^{n}=\operatorname{Ker}J\oplus\operatorname{Im}J$
    (on identifie $\R^{n}$ et $\mathcal{M}_{n,1}(\R)$), et en déduire que
    $J$ est semblable à $\operatorname{diag}(n,0,\dots,0)$.
  \item Montrer que $M(a,b)\,M(c,d)=M(ac,\ ad+bc+nbd)$. Que peut-on en
    déduire pour l'ensemble des matrices $M(a,b)$ ?
  \item Montrer que $M(a,b)$ est inversible si et seulement si $a\neq0$ et
    $a+nb\neq0$, et déterminer alors son inverse sous la forme $M(c,d)$.
  \item Pour $p\in\N$, montrer que
    $M(a,b)^{p}=a^{p}I_n+\dfrac{(a+nb)^{p}-a^{p}}{n}\,J$, et calculer
    $\det M(a,b)$.
  \item Application : pour $n=3$, calculer $(2I_3+J)^{-1}$.
\end{questions}
\end{exoEval}

\begin{exoEval}{5 points}{Records}
Soit $n\geqslant1$. On reçoit successivement $n$ candidats, tous de niveaux
différents ; leurs rangs (du moins bon, $1$, au meilleur, $n$) forment une
permutation $\sigma$ de $\ens{1,\dots,n}$, choisie selon la loi uniforme.
Pour $1\leqslant k\leqslant n$, on dit que le $k$-ième candidat est un
\emph{record} si $\sigma(k)>\sigma(i)$ pour tout $i<k$ ; on note $R_k$ cet
événement et $N_n$ le nombre de records.
\begin{questions}
  \item Montrer que $P(R_k)=\dfrac1k$.
  \item Soient $1\leqslant j<k\leqslant n$. Montrer que
    $P(R_j\cap R_k)=\dfrac{1}{jk}$. Que peut-on en déduire ?
  \item Calculer l'espérance et la variance de $N_n$ ; on notera
    $H_n=\sum_{k=1}^{n}\frac1k$.
  \item Montrer que $H_n\sim\ln n$. En déduire, à l'aide de l'inégalité de
    Bienaymé-Tchebychev, que pour tout $\varepsilon>0$,
    $P\Bigl(\Bigl|\dfrac{N_n}{H_n}-1\Bigr|\geqslant\varepsilon\Bigr)\xrightarrow[n\to+\infty]{}0$.
  \item Déterminer la loi de $N_3$.
\end{questions}
\end{exoEval}

\begin{exoEval}{5 points}{Polynômes de Bernoulli}
\begin{questions}
  \item Montrer qu'il existe une unique suite $(B_n)_{n\in\N}$ de polynômes
    réels telle que $B_0=1$ et, pour tout $n\geqslant1$,
    \[B_n'=nB_{n-1}\qquad\text{et}\qquad\int_0^{1}B_n(t)\,\mathrm{d}t=0.\]
  \item Calculer $B_1$, $B_2$, $B_3$. Montrer que $B_n$ est unitaire de
    degré $n$.
  \item Montrer que $B_n(1)=B_n(0)$ pour $n\geqslant2$, puis que
    $B_n(X+1)-B_n(X)=nX^{n-1}$ pour $n\geqslant1$.
  \item En déduire que pour $N\geqslant1$ et $p\geqslant1$ :
    $\sum_{k=0}^{N-1}k^{p}=\dfrac{B_{p+1}(N)-B_{p+1}(0)}{p+1}$. Retrouver la
    somme des carrés des $N-1$ premiers entiers.
  \item Montrer que $B_n(1-X)=(-1)^{n}B_n(X)$ pour tout $n$. En déduire que
    $B_{2p+1}\bigl(\frac12\bigr)=B_{2p+1}(0)=0$ pour $p\geqslant1$.
\end{questions}
\end{exoEval}

\begin{exoEval}{5 points}{L'équation de Cauchy}
Soit $f:\R\to\R$ telle que $f(x+y)=f(x)+f(y)$ pour tous réels $x,y$. On pose
$a=f(1)$.
\begin{questions}
  \item Montrer que $f(0)=0$, que $f$ est impaire, et que $f(nx)=nf(x)$ pour
    tous $n\in\Z$ et $x\in\R$.
  \item Montrer que $f(r)=ar$ pour tout rationnel $r$.
  \item On suppose $f$ continue. Montrer que $f(x)=ax$ pour tout réel $x$.
  \item On suppose seulement $f$ croissante. Montrer que la conclusion
    subsiste.
  \item Soit $g:\R\to\R$ continue telle que
    $g\bigl(\frac{x+y}{2}\bigr)=\frac{g(x)+g(y)}{2}$ pour tous $x,y$. Montrer
    que $g$ est affine. \emph{Considérer $x\mapsto g(x)-g(0)$.}
\end{questions}
\end{exoEval}

% !TEX root = ../main.tex
%% Sujet type Mines-Ponts — Polynômes de Legendre et méthode de Gauss.
%% Résultats vérifiés par verifications/eval_cc.py.

\sujetEval{C2}{Sujet type Mines-Ponts · Mathématiques}{Sujet type Mines-Ponts}
  {Durée : \textbf{3 heures}\\Calculatrice interdite\\Polynômes de Legendre}{4}
  {Partie A/Les polynômes de Legendre/Q1 à Q5,Partie B/Les racines de $L_n$/Q6,%
   Partie C/La méthode de Gauss/Q7 à Q11,Partie D/Une relation de récurrence/Q12}
\enParties
\consigne{Sujet inédit, conçu dans l'esprit des épreuves de la banque
Mines-Ponts : un problème long, dont les parties s'enchaînent ; un résultat
non démontré peut être admis pour la suite, à condition de le signaler.}

On munit $\R[X]$ de l'application $(P,Q)\mapsto\langle P,Q\rangle$
définie par
\[\langle P,Q\rangle=\int_{-1}^{1}P(t)\,Q(t)\,\mathrm{d}t.\]
Pour $n\in\N$, on pose
\[U_n=(X^{2}-1)^{n}\qquad\text{et}\qquad L_n=\frac{1}{2^{n}\,n!}\,U_n^{(n)}
  \qquad(\text{formule de Rodrigues}),\]
et l'on admet que
$\displaystyle\int_{-1}^{1}(1-t^{2})^{n}\,\mathrm{d}t=\frac{2^{2n+1}(n!)^{2}}{(2n+1)!}$.

\begin{exoEval}[Les polynômes de Legendre]{7 points}{}
\begin{questions}
  \item Montrer que $\langle\cdot,\cdot\rangle$ est un produit scalaire sur
    $\R[X]$.
  \item Calculer $L_0$, $L_1$, $L_2$ et $L_3$. Montrer que $L_n$ est de degré
    $n$, de coefficient dominant $c_n=\dfrac{(2n)!}{2^{n}(n!)^{2}}$, et qu'il a
    la parité de $n$.
  \item En appliquant la formule de Leibniz à $U_n=(X-1)^{n}(X+1)^{n}$,
    montrer que $L_n(1)=1$ et $L_n(-1)=(-1)^{n}$.
  \item Montrer que $U_n^{(k)}(1)=U_n^{(k)}(-1)=0$ pour $0\leqslant k\leqslant n-1$.
    En déduire, par $n$ intégrations par parties, que
    $\langle L_n,Q\rangle=0$ pour tout $Q\in\R_{n-1}[X]$, puis que
    $(L_0,\dots,L_n)$ est une base orthogonale de $\R_n[X]$.
  \item Montrer que $\lVert L_n\rVert^{2}=\langle L_n,c_nX^{n}\rangle$, puis
    que $\lVert L_n\rVert^{2}=\dfrac{2}{2n+1}$.
\end{questions}
\end{exoEval}

\begin{exoEval}[Les racines de $L_n$]{3 points}{}
\begin{questions}[start=6]
  \item Soit $n\geqslant1$. On note $x_1<\dots<x_p$ les racines de $L_n$
    situées dans $\intOO{-1}{1}$ et de multiplicité impaire ($p\geqslant0$),
    et $Q=\prod_{i=1}^{p}(X-x_i)$ (avec $Q=1$ si $p=0$). Montrer que $L_nQ$
    garde un signe constant sur $\intFF{-1}{1}$, puis que $p=n$. En déduire
    que $L_n$ possède $n$ racines simples, toutes dans $\intOO{-1}{1}$.
\end{questions}
\end{exoEval}

\begin{exoEval}[La méthode de Gauss]{7 points}{}
On fixe $n\geqslant1$ ; on note $x_1<\dots<x_n$ les racines de $L_n$,
$\ell_1,\dots,\ell_n$ les polynômes de Lagrange associés
($\ell_i(x_j)=\delta_{ij}$, $\deg\ell_i=n-1$) et $w_i=\int_{-1}^{1}\ell_i$.
\begin{questions}[start=7]
  \item Montrer que $\displaystyle\int_{-1}^{1}P=\sum_{i=1}^{n}w_iP(x_i)$ pour
    tout $P\in\R_{n-1}[X]$.
  \item En divisant $P$ par $L_n$, montrer que cette formule reste exacte
    pour tout $P\in\R_{2n-1}[X]$.
  \item Montrer qu'elle est fausse pour $P=\prod_{i=1}^{n}(X-x_i)^{2}$. Montrer
    que les poids $w_i$ sont strictement positifs, de somme $2$.
  \item Déterminer les nœuds et les poids pour $n=2$ et pour $n=3$.
  \item On donne $2\operatorname{sh}1\approx2{,}3504$,
    $\mathrm{e}^{1/\sqrt3}+\mathrm{e}^{-1/\sqrt3}\approx2{,}3427$ et
    $\mathrm{e}^{\sqrt{0{,}6}}+\mathrm{e}^{-\sqrt{0{,}6}}\approx2{,}6306$.
    Comparer à $\int_{-1}^{1}\mathrm{e}^{t}\,\mathrm{d}t$ les approximations
    fournies par la méthode de Gauss pour $n=2$ et $n=3$, et par la méthode
    des trapèzes à trois points $-1,0,1$.
\end{questions}
\end{exoEval}

\begin{exoEval}[Une relation de récurrence]{3 points}{}
\begin{questions}[start=12]
  \item Soit $n\geqslant1$. Montrer que $XL_n$ est combinaison linéaire de
    $L_{n+1}$, $L_n$ et $L_{n-1}$, puis que
    \[(n+1)\,L_{n+1}=(2n+1)\,X\,L_n-n\,L_{n-1}.\]
    Retrouver ainsi $L_4$.
\end{questions}
\end{exoEval}

% !TEX root = ../main.tex
%% Sujet type Centrale-Supélec — Chaînes de Markov et classement PageRank.
%% Résultats vérifiés par verifications/eval_cc.py (y compris les programmes).

\sujetEval{C3}{Sujet type Centrale-Supélec · Mathématiques}{Sujet type Centrale-Supélec}
  {Durée : \textbf{4 heures}\\Calculatrice interdite\\Markov, PageRank, Python}{4}
  {Partie A/Généralités/Q1 et Q2,Partie B/Une chaîne à deux états/Q3 à Q5,%
   Partie C/Le théorème de convergence/Q6 à Q8,Partie D/Le classement PageRank/Q9 à Q12}
\enParties
\consigne{Sujet inédit, conçu dans l'esprit des épreuves de la banque
Centrale-Supélec : un problème de modélisation, qui mêle algèbre linéaire,
probabilités, analyse et programmation en Python.}

Pour classer $n$ pages web, on imagine un internaute qui, à chaque instant,
quitte la page où il se trouve en suivant un lien. On note $X_k$ la page
visitée à l'instant $k$, et l'on suppose que
$P(X_{k+1}=j\mid X_k=i)=p_{ij}$ ne dépend pas de $k$. La matrice
$P=(p_{ij})$ est \emph{stochastique} : ses coefficients sont positifs et la
somme de chacune de ses lignes vaut $1$. Un \emph{vecteur de probabilité}
est une ligne $\mu=(\mu_1,\dots,\mu_n)$ de réels positifs de somme $1$ ; on
note $\mu_k$ la loi de $X_k$, et $\lVert x\rVert_1=\sum_j|x_j|$ pour toute
ligne $x$. Un vecteur de probabilité $\pi$ tel que $\pi P=\pi$ est dit
\emph{stationnaire}.

\begin{exoEval}[Généralités]{3 points}{}
\begin{questions}
  \item À l'aide de la formule des probabilités totales, montrer que
    $\mu_{k+1}=\mu_kP$, puis que $\mu_k=\mu_0P^{k}$.
  \item Montrer qu'un produit de matrices stochastiques est stochastique,
    et que $\mu P$ est un vecteur de probabilité si $\mu$ en est un.
\end{questions}
\end{exoEval}

\begin{exoEval}[Une chaîne à deux états]{4 points}{}
Soient $a,b\in\intFF01$ avec $a+b>0$, et
\[P=\begin{pmatrix}1-a&a\\b&1-b\end{pmatrix},\qquad
  R=\frac{1}{a+b}\begin{pmatrix}b&a\\b&a\end{pmatrix},\qquad
  S=\frac{1}{a+b}\begin{pmatrix}a&-a\\-b&b\end{pmatrix}.\]
\begin{questions}[start=3]
  \item Vérifier que $R+S=I_2$, $R^{2}=R$, $S^{2}=S$, $RS=SR=0$, et que
    $P=R+(1-a-b)\,S$.
  \item En déduire $P^{k}$ pour $k\in\N$. Montrer que si
    $(a,b)\neq(1,1)$, la suite $(\mu_k)$ converge, quel que soit $\mu_0$,
    vers $\pi=\bigl(\frac{b}{a+b},\frac{a}{a+b}\bigr)$, et que $\pi$ est
    stationnaire. Que se passe-t-il si $a=b=1$ ?
  \item Un jour ensoleillé est suivi d'un jour de pluie avec probabilité
    $0{,}3$, un jour de pluie d'un jour ensoleillé avec probabilité $0{,}4$.
    Quelle est, à long terme, la proportion de jours ensoleillés ?
\end{questions}
\end{exoEval}

\begin{exoEval}[Le théorème de convergence]{6 points}{}
Soit $Q$ une matrice stochastique de taille $n$ dont tous les coefficients
sont supérieurs ou égaux à un réel $\delta>0$.
\begin{questions}[start=6]
  \item Montrer que $n\delta\leqslant1$. Soit $x$ une ligne telle que
    $\sum_ix_i=0$ ; montrer que $(xQ)_j=\sum_ix_i\,(q_{ij}-\delta)$, puis que
    $\lVert xQ\rVert_1\leqslant(1-n\delta)\,\lVert x\rVert_1$.
  \item Soit $\mu_0$ un vecteur de probabilité et $\mu_k=\mu_0Q^{k}$.
    Montrer que
    $\lVert\mu_{k+1}-\mu_k\rVert_1\leqslant(1-n\delta)^{k}\lVert\mu_1-\mu_0\rVert_1$.
    En déduire que chaque suite de coordonnées de $(\mu_k)$ converge, et que
    la limite $\pi$ est un vecteur de probabilité stationnaire pour $Q$.
  \item Montrer que $Q$ admet un unique vecteur de probabilité stationnaire,
    et que $\lVert\mu_k-\pi\rVert_1\leqslant2(1-n\delta)^{k}$.
\end{questions}
\end{exoEval}

\begin{exoEval}[Le classement PageRank]{7 points}{}
\noindent\begin{minipage}[c]{0.62\linewidth}
On considère le réseau de quatre pages ci-contre : une flèche de $i$ vers
$j$ représente un lien de la page $i$ vers la page $j$. L'internaute suit,
de façon équiprobable, l'un des liens de la page où il se trouve.
\end{minipage}\hfill
\begin{minipage}[c]{0.34\linewidth}\centering
\begin{tikzpicture}[scale=0.8,line cap=round,>={Stealth[length=2.2mm]},
    pg/.style={circle,draw=otNavy,line width=0.9pt,fill=otNavyTint,minimum size=7.5mm,
      font=\otSemi\fontsize{9}{10}\selectfont,text=otNavy}]
  \node[pg] (p1) at (0,1.9) {1};
  \node[pg] (p2) at (3,1.9) {2};
  \node[pg] (p3) at (3,0) {3};
  \node[pg] (p4) at (0,0) {4};
  \draw[->,otInk,line width=0.7pt] (p1) -- (p2);
  \draw[->,otInk,line width=0.7pt] (p2) -- (p3);
  \draw[->,otInk,line width=0.7pt] (p1) to[bend left=12] (p3);
  \draw[->,otInk,line width=0.7pt] (p3) to[bend left=12] (p1);
  \draw[->,otInk,line width=0.7pt] (p4) -- (p1);
  \draw[->,otInk,line width=0.7pt] (p4) -- (p3);
\end{tikzpicture}
\end{minipage}\par\vspace{1mm}
\begin{questions}[start=9]
  \item Écrire $P$ et vérifier que $\bigl(\frac25,\frac15,\frac25,0\bigr)$ est
    stationnaire. La partie C s'applique-t-elle à $P$ ? Pourquoi un score
    nul pour la page $4$ est-il insatisfaisant ?
  \item Soit $\alpha\in\intOO01$ et $G=\alpha P+\frac{1-\alpha}{n}J$, où $J$
    est la matrice dont tous les coefficients valent $1$. Interpréter $G$.
    Montrer que $G$ admet un unique vecteur de probabilité stationnaire
    $\pi$ (le \emph{PageRank}) et que
    $\lVert\mu_k-\pi\rVert_1\leqslant2\alpha^{k}$.
  \item Sans calculer $\pi$, montrer que $\pi_4=\frac{1-\alpha}{4}$.
  \item On considère la fonction Python suivante, où \texttt{P} est une
    liste de listes représentant $P$.
\begin{python}
def pagerank(P, alpha, eps):
    n = len(P)
    mu = [1 / n] * n
    while True:
        nu = [(1 - alpha) / n
              + alpha * sum(mu[i] * P[i][j] for i in range(n))
              for j in range(n)]
        if sum(abs(nu[j] - mu[j]) for j in range(n)) < eps:
            return nu
        mu = nu
\end{python}
    \begin{enumerate}[label=\alph*)]
      \item Montrer que la liste \texttt{nu} représente $\mu G$, et que la
        boucle se termine. Majorer le nombre de passages dans la boucle pour
        $\alpha=0{,}85$ et $\varepsilon=10^{-6}$ (on donne
        $\ln(2\cdot10^{6})\approx14{,}51$ et $\ln0{,}85\approx-0{,}1625$).
      \item Avec $\alpha=0{,}85$, le programme renvoie
        $(0{,}3797;\ 0{,}1989;\ 0{,}3839;\ 0{,}0375)$. Classer les pages ;
        commenter.
      \item Évaluer le coût d'un passage dans la boucle. Écrire une fonction
        \texttt{pagerank\_creux(succ, alpha, eps)}, où \texttt{succ[i]} est
        la liste (non vide) des pages liées depuis la page \texttt{i}, dont
        chaque passage coûte $O(n+\text{nombre de liens})$ opérations.
    \end{enumerate}
\end{questions}
\end{exoEval}

% !TEX root = ../main.tex
%% Sujet type X-ENS — Inégalités de Hardy et de Carleman.
%% Résultats vérifiés par verifications/eval_cc.py.

\sujetEval{C4}{Sujet type X-ENS · Mathématiques}{Sujet type X-ENS}
  {Durée : \textbf{4 heures}\\Calculatrice interdite\\Inégalités de moyennes}{4}
  {Partie A/L'inégalité de Hardy/Q1 à Q4,Partie B/La constante 4 est optimale/Q5 et Q6,%
   Partie C/L'inégalité de Carleman/Q7 à Q10,Partie D/La constante $\mathrm{e}$ est optimale/Q11 et Q12}
\enParties
\consigne{Sujet inédit, conçu dans l'esprit des épreuves de l'X et des ENS :
peu de calculs, mais des raisonnements à construire. On pourra utiliser
l'inégalité arithmético-géométrique (devoir maison \numero1) et l'inégalité
$\bigl(1+\frac1k\bigr)^{k}<\mathrm{e}$, valable pour $k\geqslant1$.}

Dans tout le problème, $(a_n)_{n\geqslant1}$ est une suite de réels
positifs, et l'on pose, pour $n\geqslant1$,
\[A_n=\frac{a_1+a_2+\dots+a_n}{n}\qquad\text{et}\qquad G_n=\bigl(a_1a_2\cdots a_n\bigr)^{1/n},\]
ainsi que $A_0=0$. On compare les séries $\sum A_n^{2}$ et $\sum a_n^{2}$,
puis $\sum G_n$ et $\sum a_n$.

\begin{exoEval}[L'inégalité de Hardy]{6 points}{}
\begin{questions}
  \item Montrer que $a_n=nA_n-(n-1)A_{n-1}$ pour tout $n\geqslant1$.
  \item En déduire que
    $A_n^{2}-2a_nA_n\leqslant(n-1)A_{n-1}^{2}-nA_n^{2}$ pour $n\geqslant1$.
  \item Soit $N\geqslant1$. Montrer que
    $\sum_{n=1}^{N}A_n^{2}\leqslant2\sum_{n=1}^{N}a_nA_n$, puis que
    \[\sum_{n=1}^{N}A_n^{2}\leqslant4\sum_{n=1}^{N}a_n^{2}.\]
  \item En déduire que si $\sum a_n^{2}$ converge, alors $\sum A_n^{2}$
    converge et $\sum_{n=1}^{+\infty}A_n^{2}\leqslant4\sum_{n=1}^{+\infty}a_n^{2}$
    (\emph{inégalité de Hardy}).
\end{questions}
\end{exoEval}

\begin{exoEval}[La constante 4 est optimale]{4 points}{}
On fixe $N\geqslant1$ et l'on choisit $a_n=\frac{1}{\sqrt n}$ pour
$n\leqslant N$, $a_n=0$ pour $n>N$.
\begin{questions}[start=5]
  \item Montrer que pour $n\geqslant1$ :
    $2\sqrt{n+1}-2\leqslant\sum_{k=1}^{n}\dfrac{1}{\sqrt k}\leqslant2\sqrt n-1$.
  \item En déduire qu'il existe une constante $K$, indépendante de $N$,
    telle que $\sum_{n=1}^{N}A_n^{2}\geqslant4H_N-K$, où
    $H_N=\sum_{n=1}^{N}\frac1n$. Montrer que $4$ ne peut être remplacé, dans
    l'inégalité de Hardy, par une constante strictement plus petite.
\end{questions}
\end{exoEval}

\begin{exoEval}[L'inégalité de Carleman]{6 points}{}
Pour $k\geqslant1$, on pose $c_k=\dfrac{(k+1)^{k}}{k^{k-1}}$.
\begin{questions}[start=7]
  \item Montrer que $c_1c_2\cdots c_n=(n+1)^{n}$ pour tout $n\geqslant1$.
  \item En appliquant l'inégalité arithmético-géométrique aux réels
    $a_1c_1,\dots,a_nc_n$, montrer que
    $G_n\leqslant\dfrac{1}{n(n+1)}\sum_{k=1}^{n}a_kc_k$.
  \item Montrer que $\sum_{n=k}^{N}\dfrac{1}{n(n+1)}\leqslant\dfrac1k$, puis que
    \[\sum_{n=1}^{N}G_n\leqslant\sum_{k=1}^{N}a_k\Bigl(1+\frac1k\Bigr)^{k}.\]
  \item En déduire que si $\sum a_n$ converge, alors $\sum G_n$ converge et
    $\sum_{n=1}^{+\infty}G_n\leqslant\mathrm{e}\sum_{n=1}^{+\infty}a_n$
    (\emph{inégalité de Carleman}).
\end{questions}
\end{exoEval}

\begin{exoEval}[La constante $\mathrm{e}$ est optimale]{4 points}{}
\begin{questions}[start=11]
  \item Montrer que pour $n\geqslant1$ :
    $n\ln n-n+1\leqslant\ln n!\leqslant(n+1)\ln(n+1)-n$. En déduire que
    $(n!)^{1/n}\sim\dfrac{n}{\mathrm{e}}$.
  \item On fixe $N\geqslant1$ et l'on choisit $a_n=\frac1n$ pour
    $n\leqslant N$, $a_n=0$ pour $n>N$. Soit $\varepsilon\in\intOO01$.
    Montrer qu'il existe $n_0$, indépendant de $N$, tel que
    $\sum_{n=1}^{N}G_n\geqslant(1-\varepsilon)\,\mathrm{e}\,(H_N-H_{n_0})$
    pour $N\geqslant n_0$. Conclure.
\end{questions}
\end{exoEval}

